\par\Needspace{6\baselineskip}\medskip
\begin{center}\hypertarget{native-197-article147}{}\pdfbookmark[2]{147.}{article-147}147.\end{center}
Streng genommen steht es uns auch hier, wie bei den zehn oben dargelegten Verfahren, frei, welche neuen Werte für $x$ und $y$ wir in der zweiten Hypothese ansetzen, sofern sie den oben erläuterten allgemeinen Bedingungen genügen. Es ist jedoch offensichtlich von großem Vorteil, gleich von genaueren Werten ausgehen zu können. Daher wäre es bei diesem Verfahren unklug, die zweiten Werte aufs Geratewohl zu wählen. Aus der Natur der Sache ist leicht einzusehen, dass bei kleinen Fehlern der ersten Werte von $P,Q$ die Größen $P',Q'$ wesentlich genauere Werte liefern, sofern die heliozentrische Bewegung mäßig ist. Wir werden daher stets $P',Q'$ als zweite Werte von $P,Q$ übernehmen, oder $\log P',\log Q'$ als zweite Werte von $x,y$, falls $\log P,\log Q$ zur Bezeichnung der ersten Werte angesetzt wurden.

Bei dieser zweiten Hypothese bleiben sämtliche in den Formeln 1—20 angegebenen Vorarbeiten unverändert, und die Rechnung wird auf genau dieselbe Weise wiederholt. Zuerst bestimmt man den Winkel $\omega$, dann $z,r',\dfrac{n'r'}{n}\mathpunct{\raisebox{.45ex}{,}}$ $\dfrac{n'r'}{n''}$
$\zeta,r,\zeta'',r'',f',f,f''$. Je nachdem, wie stark sich die neuen Werte dieser Größen von den ersten unterscheiden, lässt sich leicht beurteilen, ob sich auch eine erneute Berechnung der Aberrationskorrektur der Zeiten lohnt. Andernfalls bleiben die Zeitintervalle und damit auch die Größen $\theta$ und $\theta''$ wie zuvor. Schließlich gewinnt man aus $f,r',r''$; $f'',r,r'$ und den Zeitintervallen die Größen $\eta,\eta''$ und daraus neue Werte von $P',Q'$. Diese werden sich von den Werten aus der ersten Hypothese meist weit weniger unterscheiden als jene von den ersten Werten von $P,Q$. Die zweiten Werte von $X,Y$ werden also viel
\SourceLeaf{198}kleiner sein als die ersten. Die zweiten Werte von $P',Q'$ werden als dritte Werte von $P,Q$ übernommen, und mit ihnen wird die Rechnung erneut durchgeführt.
Wie sich auf diese Weise aus der zweiten Hypothese genauere Zahlen ergaben als aus der ersten, so werden sich aus der dritten wiederum genauere ergeben als aus der zweiten. Man könnte die dritten Werte von $P',Q'$ als vierte Werte von $P,Q$ übernehmen und die Rechnungen so lange wiederholen, bis man zu einer Hypothese gelangt, in der $X$ und $Y$ als verschwindend betrachtet werden können. Scheint die dritte Hypothese noch nicht auszureichen, so ist es zweckmäßiger, die für die vierte Hypothese anzusetzenden Werte von $P,Q$ nach dem in Art. 120, 121 erläuterten Verfahren aus den ersten drei Hypothesen abzuleiten. So erhält man die Näherung schneller und wird nur selten zu einer fünften Hypothese übergehen müssen.
